\section{Comparación de métodos y análisis de rendimiento}

\subsection{MDLM frente a BD3LM y SoLM}

BD3LM (Modelo de lenguaje de difusión por bloques) es otro método de modelo de lenguaje de difusión que admite caché KV. Su enfoque consiste en dividir secuencias largas en bloques de tamaño fijo (por ejemplo, 100 tokens por bloque), realizar el proceso de desruido mediante difusión dentro de cada bloque y avanzar entre bloques utilizando un modo autorregresivo.

\begin{table}[H]
\centering
\caption{Comparación de tres métodos de modelo de lenguaje de difusión}
\begin{tabular}{@{}lccc@{}}
\toprule
\textbf{Característica} & \textbf{MDLM} & \textbf{BD3LM} & \textbf{SoLM} \\
\midrule
Tipo de atención & Bidireccional & Bidireccional dentro del bloque + causal entre bloques & Causal (tras ordenar) \\
Caché KV & No admite admitida parcialmente (solo entre bloques) & Admitida completamente \\
Generación paralela & Paralelismo global & Paralelismo dentro del bloque & Paralelismo global \\
Velocidad de inferencia & Lenta & Moderada & Rápida \\
Calidad de generación & Alta & Media-alta & Alta \\
\bottomrule
\end{tabular}
\end{table}

\begin{warningbox}{Degradación de calidad de BD3LM en muestreo con pocos pasos}
BD3LM mantiene una buena calidad al aumentar el número de pasos de muestreo, pero su calidad se degrada significativamente cuando se reduce drásticamente el número de pasos de muestreo (es decir, generar más palabras en paralelo por paso), e incluso puede ser inferior a la de MDLM. Esto se debe a que, tras reducir los pasos de difusión dentro del bloque, el efecto de desruido es insuficiente. En comparación, SoLM mantiene una buena calidad incluso con muestreo reducido.
\end{warningbox}

\subsection{Curva de compensación entre velocidad y calidad}

Sobre la frontera de Pareto entre perplexidad y tiempo de muestreo:

\begin{itemize}
    \item \textbf{MDLM}: Perplexidad óptima, pero velocidad de inferencia más lenta (la curva se sitúa en la parte inferior derecha)
    \item \textbf{BD3LM}: Es más rápido que MDLM con un alto número de pasos de muestreo, pero la calidad disminuye drásticamente tras reducir los pasos
    \item \textbf{SoLM}: La frontera de Pareto es \textbf{estrictamente superior} a la de MDLM y BD3LM; es decir, ofrece mayor calidad para cualquier velocidad dada o mayor velocidad para cualquier calidad dada
\end{itemize}

En pruebas basadas en 110 millones de parámetros y un contexto de longitud 1024 tokens, el tiempo de muestreo de SoLM logró una reducción significativa en comparación con MDLM.

\subsection{Resumen de la sección}

SoLM es estrictamente superior a MDLM y BD3LM en términos de compensación entre velocidad y calidad, lo cual se debe principalmente a su soporte completo para caché KV. Aunque el diseño por bloques de BD3LM también admite parcialmente la caché KV, sigue estando limitado por los cuellos de botella de la atención bidireccional dentro del bloque.

\newpage

